\chapter{From SPARQL to a Plan}
\label{ch:plans}

A query states \emph{what} to find. The engine must decide \emph{how}. This
chapter follows a SPARQL query from text to an executable tree of operators,
and introduces the iterator model that runs it.

\section{Parsing and binding}

The \term{parser} turns text into an abstract syntax tree: prefixes, the
projected variables, the triple patterns, the filters. Syntax that the engine
does not support should fail here, with a line and a column, rather than half
work later.

The \term{binder} then resolves names against the database:

\begin{itemize}
\item Prefixed names become full IRIs: \code{ub:Student} becomes\\ \code{<http://example.edu/univ\#Student>}.
\item Every constant becomes a dictionary id. For a query, a constant that is not in the dictionary can match nothing; for an insert, it is added.
\item Every variable gets a \term{slot}: a column number, assigned in order of first appearance.
\end{itemize}

After binding, a row is simply an array of ids indexed by slot. The same
variable in two patterns is the same slot, and that shared slot is the join key.

\section{Logical and physical plans}

A \term{logical plan} says what to compute in terms of algebra: scan each
pattern, join the results, filter, project. A \term{physical plan} picks an
algorithm for each step: which index and range to scan, which join algorithm,
in what order. The planner produces a correct physical plan; the optimizer
(\cref{ch:joins}) rewrites it into a fast one that returns the same rows.

For the query

\begin{lstlisting}
SELECT ?name WHERE {
  ?s a ub:Student . ?s ub:name ?name . ?s ub:age ?age .
  FILTER (?age > 21)
}
\end{lstlisting}

a straightforward plan is left-deep, joining the patterns in the order written:

\begin{lstlisting}
Projection [?name]
  Filter (?age > 21)
    NestedLoopJoin
      NestedLoopJoin
        TripleScan {?s rdf:type ub:Student}
        TripleScan {?s ub:name ?name}
      TripleScan {?s ub:age ?age}
\end{lstlisting}

\section{The iterator (Volcano) model}

Graefe's \term{Volcano} model \cite{graefe} gives every operator the same
three-call interface:

\begin{itemize}
\item \code{Init()}: prepare, or restart from the beginning.
\item \code{Next(\&row)}: produce the next row and return true, or return false when exhausted.
\item destruction: release everything held (for a scan, its pinned leaf).
\end{itemize}

The root operator calls \code{Next} on its child, which calls \code{Next} on
its children, and so on down to the scans. Rows flow up \emph{one at a time},
on demand. Three properties make this model the default for decades:

\begin{description}[leftmargin=1.2em]
\item[Composability.] Any operator can sit above any other, because they share an interface.
\item[Pipelining.] A row can travel from a scan to the client without any intermediate result being materialized. \code{LIMIT 10} stops the whole tree after ten rows.
\item[Bounded memory.] Only operators that must see all their input first (sort, hash-join build, aggregation) hold rows.
\end{description}

Its cost is one virtual call per row per operator. Modern analytical engines
reduce that with \term{vectorized} execution (batches of rows per call) or by
compiling the plan into tight loops, but the logical structure is unchanged.

\begin{keyidea}
Planning builds a tree; execution pulls rows through it. Keeping the two
separate is what allows \code{\textbackslash explain} to print a plan without
running it, and an optimizer to swap one tree for another.
\end{keyidea}

\section{The operators}

\begin{description}[leftmargin=1.2em]
\item[Triple scan.] Reads the triples that match one pattern. With an index, the bound positions choose a permutation and a prefix range (\cref{ch:rdf-indexing}). If a variable occurs twice in one pattern (\code{?x ex:knows ?x}), the scan also checks that those positions are equal.
\item[Join.] Combines rows from two children whose shared slots agree (\cref{ch:joins}).
\item[Filter.] Passes rows whose expression is true. In SPARQL a type error makes the expression false, not the query fail.
\item[Projection.] Keeps the selected slots.
\item[Distinct, limit, offset, sort.] Distinct applies to \emph{projected} rows (two triples can share a name). Sort compares values, not ids: ids are assigned in load order, which is not value order.
\end{description}

\begin{ontodb}[: planning and execution]
OntoDB's binder assigns slots in first-appearance order and represents ``not yet
bound'' with a reserved id. The planner always produces the left-deep tree in
textual order with nested-loop joins, which is correct but rarely fast;
\code{\textbackslash explain} prints it. An executor factory maps each plan node
to its executor, and the engine pulls rows with \code{Init}/\code{Next}.
An empty pattern, with no triple to match, returns one empty row rather than
none: the identity input for a join that has nothing to bind yet.
\end{ontodb}

\begin{decision}{Unsupported syntax fails at parse time}
\dpart{Choice} \code{OPTIONAL}, \code{UNION}, property paths, aggregates, and
\code{DELETE/INSERT WHERE} are recognized by the parser and rejected there with
a line and column.
\dpart{Alternatives} Accept the full grammar and fail in the executor; or skip
what is not understood.
\dpart{Why} A query either runs completely or fails before anything happens.
Ignoring an \code{OPTIONAL} would silently change the answer, which is worse
than refusing.
\dpart{Cost} The parser must know the shape of features the engine does not
implement.
\end{decision}

\begin{decision}{Volcano, one row at a time}
\dpart{Choice} Every operator implements \code{Init} and \code{Next} and passes
one row of ids per call.
\dpart{Alternatives} Vectorized execution (a batch of rows per call), or
compiling the plan into loops.
\dpart{Why} It is the simplest model that composes, it makes \code{LIMIT} stop
the whole tree, and it keeps pin ownership obvious: a scan holds one leaf while
it has a current row.
\dpart{Cost} A virtual call per row per operator. Large analytical scans would
be several times faster vectorized.
\end{decision}

\begin{decision}{Stub plan nodes before their executors}
\dpart{Choice} \code{Sort}, \code{Distinct}, \code{Limit}, \code{Insert}, and
\code{Delete} exist as plan nodes that \code{\textbackslash explain} prints
before their executors are written.
\dpart{Alternatives} Add a plan node only together with its executor.
\dpart{Why} The plan's shape and the planner's tests can be fixed first; each
executor then arrives into a slot that is already wired and printed.
\dpart{Cost} For a while, \code{\textbackslash explain} shows plans that cannot
run, and running one throws ``not built yet''.
\end{decision}

\begin{decision}{\code{DISTINCT} keeps a hash set}
\dpart{Choice} The distinct operator passes a row the first time it sees it,
in its child's order, and remembers every row it has passed in a hash set.
\dpart{Alternatives} Sort the rows and drop adjacent duplicates.
\dpart{Why} Hashing streams: the first answer comes out as soon as it is
found, \code{LIMIT} above it still stops early, and the order the child chose
is kept.
\dpart{Cost} Memory grows with the number of distinct rows. A sort-based
distinct could spill to disk; this one cannot.
\end{decision}

\begin{decision}{Rows have no order without \code{ORDER BY}}
\dpart{Choice} The engine promises no order unless the query has
\code{ORDER BY}. Tests compare answers as multisets.
\dpart{Alternatives} Always return rows in a fixed order, for example sorted by
id.
\dpart{Why} The order a plan produces depends on the store, the indexes, and
the join operator: the memory store returns triples in insertion order, the
indexed store in key order. Promising an order would force a sort on every
query.
\dpart{Cost} \code{LIMIT} without \code{ORDER BY} may return different rows on
the two backends, or under different join operators. Both answers are correct,
which is why the differential test checks such queries only when the limit is
larger than the answer.
\end{decision}

\begin{status}
Parsing, binding, planning, triple scans, the three joins, filter,
projection, distinct, limit, and offset run on both backends (plans 4.1 to
4.4). Sort is parsed and planned; its executor is plan 4.5.
\end{status}

\begin{checkyourself}
\item In \code{?s a ub:Student . ?s ub:name ?name . ?s ub:age ?age}, which slot does each variable get?
\item Why does \code{DISTINCT} apply to projected rows, not to the rows the scans produce?
\item Why can't \code{ORDER BY} compare term ids?
\item A plan is \code{Limit 10} over a join of two large scans. How much work does the Volcano model do, and how much would a model that materializes each operator's output do?
\item Why does an empty pattern produce one empty row rather than none?
\item \code{SELECT ?name WHERE \{ ... \} LIMIT 2} returns Alice and Bob on the memory backend and Alice and Cara on the indexed one. Which is wrong?
\end{checkyourself}

\section*{Further reading}
Graefe \cite{graefe}; the W3C SPARQL 1.1 recommendation \cite{sparql}.
